\section*{Chapter \ref{ch:g7:lists} --- Lists}

\begin{solution}{exo:g7:lists:1}
\code{L[0]} is 4, \code{L[2]} is 15 (the third \omterm{def:g7:lists:list}{item}), \code{L[-1]} is 42 (the
last \omterm{def:g7:lists:list}{item}) and \code{len(L)} is 6.
\end{solution}

\begin{solution}{exo:g7:lists:2}
\code{[1, 2, 3, 4]} after \code{append}, \code{[2, 3, 4]} after \code{pop(0)},
which removes the \omterm{def:g7:lists:list}{item} at \omterm{def:g7:lists:list}{index} 0, and \code{[2, 9, 3, 4]} after the
\code{insert}. It prints \code{[2, 9, 3, 4]}.
\end{solution}

\begin{solution}{exo:g7:lists:3}
6, 2 and 8, one per line.
\end{solution}

\begin{solution}{exo:g7:lists:4}
\omcode{code/grade-7/02-lists/exercises/average.py}{1}{6}{The average daily
rainfall.}\label{lst:g7:lists:sol-average}

It prints \code{average: 1.82 mm a day}: 56.3 millimetres in 31 days.
\end{solution}

\begin{solution}{exo:g7:lists:5}
\code{range(len(temps) + 1)} counts 0, 1, 2, 3, 4, but the last \omterm{def:g7:lists:list}{index} of a \omterm{def:g7:lists:list}{list}
of 4 \omterm{def:g7:lists:list}{items} is 3: \code{temps[4]} is an \code{IndexError}. The walk must use
\code{range(len(temps))}, or simply \code{for t in temps:}.
\end{solution}

\begin{solution}{exo:g7:lists:6}
\code{largest} starts at 0, which is higher than every temperature of the \omterm{def:g7:lists:list}{list}:
no \omterm{def:g7:lists:list}{item} is ever larger, and 0 stays. But 0 is not in the \omterm{def:g7:lists:list}{list}! The best so far
must start with an \omterm{def:g7:lists:list}{item} of the \omterm{def:g7:lists:list}{list}:

\omcode{code/grade-7/02-lists/exercises/largest_fixed.py}{1}{6}{Starting with the
first \omterm{def:g7:lists:list}{item}.}\label{lst:g7:lists:sol-largest}

It prints $-1$.
\end{solution}

\begin{solution}{exo:g7:lists:7}
\omcode{code/grade-7/02-lists/exercises/rainy_days.py}{1}{7}{Counting the rainy
days.}\label{lst:g7:lists:sol-rainy}

It prints 13, and $13 + 18 = 31$: each day is either dry or rainy.
\end{solution}

\begin{solution}{exo:g7:lists:8}
\omcode{code/grade-7/02-lists/exercises/above5.py}{1}{7}{A new \omterm{def:g7:lists:list}{list} of the heavy
days.}\label{lst:g7:lists:sol-above5}

It prints \code{[6.9, 16.1, 5.9, 7.7, 5.4]}.
\end{solution}

\begin{solution}{exo:g7:lists:9}
31 \omterm{def:g6:if-else:comparison}{comparisons}: the walk goes through all 31 \omterm{def:g7:lists:list}{items}, the first one included,
which it compares with itself. \cref{lst:g7:lists:position-of-max} starts its
walk at \omterm{def:g7:lists:list}{index} 1, since the \omterm{def:g7:lists:list}{item} at \omterm{def:g7:lists:list}{index} 0 is already the best so far: 30
\omterm{def:g6:if-else:comparison}{comparisons}. For $n$ \omterm{def:g7:lists:list}{items}, $n - 1$ \omterm{def:g6:if-else:comparison}{comparisons} are enough.
\end{solution}

\begin{solution}{exo:g7:lists:10}
Each \omterm{def:g7:lists:list}{item} is inserted at the front, so the last one read ends up first:

\omcode{code/grade-7/02-lists/exercises/reverse.py}{1}{5}{Reversing a
\omterm{def:g7:lists:list}{list}.}\label{lst:g7:lists:sol-reverse}

It prints \code{[0.3, 0.0, 6.9, 0.0, 0.7, 0.0, 0.0]}.
\end{solution}

\begin{solution}{exo:g7:lists:11}
Before the walk, \code{largest} is the first \omterm{def:g7:lists:list}{item}: the largest of the \omterm{def:g7:lists:list}{items}
seen so far, since only one has been seen. At each turn, if the new \omterm{def:g7:lists:list}{item} is
larger, it replaces \code{largest}; if not, \code{largest} stays, and it is
still at least as large as everything seen. So after each turn, \code{largest}
is the largest of the \omterm{def:g7:lists:list}{items} seen so far. After the last turn, every \omterm{def:g7:lists:list}{item} has
been seen: \code{largest} is the largest of the whole \omterm{def:g7:lists:list}{list}.
\end{solution}

\begin{solution}{exo:g7:lists:12}
The walk compares each \omterm{def:g7:lists:list}{item} with the one before it: \code{run} counts the
equal neighbours up to the current \omterm{def:g7:lists:list}{item} and \code{longest} keeps the record.
An empty \omterm{def:g7:lists:list}{list} has no run at all: it is the special case tested first, and the
fourth check.

\omcode{code/grade-7/02-lists/exercises/longest_equal.py}{1}{19}{The longest
run of equal neighbours.}\label{lst:g7:lists:sol-longest-equal}
\end{solution}

\begin{solution}{pb:g7:lists:1}
\textbf{1.} The twelfth of January, with 16.1 millimetres.

\textbf{2.} 18 dry days (the \omterm{def:g7:lists:list}{items} equal to 0.0).

\textbf{3.} From the 6th to the 9th: 4 dry days in a row. The other runs are
shorter: 2 days (1st and 2nd, 13th and 14th, 20th and 21st, 28th and 29th),
or 1 day.

\textbf{4.} Walk through the days, keeping two numbers in mind: the \omterm{def:g7:lists:list}{length} of
the dry run that ends today (one more than yesterday's if today is dry, 0 if
it rained) and the longest run met so far, with the day it ended.

\textbf{5.} \code{run} is the number of dry days in a row ending with the day
just read. A rainy day breaks the run, so the count starts again from 0.

\textbf{6.} When \code{run} beats the record, the current day is the last day
of the new longest run. The run has \code{longest} days and ends at \omterm{def:g7:lists:list}{index}
\code{longest_end}, so it starts at \code{longest_end - longest + 1}.

\textbf{7.} It agrees: 4 days from the 6th of January.

\omoutput{code/grade-7/02-lists/outputs/dry_spell.txt}{January's longest dry
run.}\label{out:g7:lists:dry-spell}

\textbf{8.} The first one: the record is beaten only by a strictly longer run
(\code{run > longest}), so a run of the same \omterm{def:g7:lists:list}{length} does not replace it.

\textbf{9.} \omterm{def:g7:lists:list}{Index} 0 gives 1 January. For 31, the \omterm{def:g4:algorithms-repeat-and-decide:loop}{loop} takes away January's 31
days: 1 February. For 58: $58 - 31 = 27$, less than 28, so 28 February; for 59:
$59 - 31 - 28 = 0$, so 1 March.

\textbf{10.} The \omterm{def:g4:programs-and-operating-system:program}{program} prints the year's result:

\omoutput{code/grade-7/02-lists/outputs/dry_year.txt}{The longest dry spell of
2025 at Hilltop.}\label{out:g7:lists:dry-year}

\textbf{11.} Yes: the year is one \omterm{def:g7:lists:list}{list}, so a run simply goes on from the last
day of a month to the first day of the next. With twelve \omterm{def:g7:lists:list}{lists}, a walk would
have to carry \code{run} from one \omterm{def:g7:lists:list}{list} into the next, and a run that crosses
two months would be easy to cut in two by mistake.

\textbf{12.} \textbf{Result.} The longest dry spell of 2025 at Hilltop lasted
11 days, from 10 August to 20 August, as \cref{out:g7:lists:dry-year} shows.
\end{solution}
